\section{CORAL 子域对齐}

\begin{frame}[t]{子域对齐：分别匹配同一类别的两域分布}
\small
第 $c$ 类子域为源域与目标域中属于该类别的样本；目标类别由模型概率估计。
\[
\text{全局对齐：}\mathcal D(P_s,P_t)
\qquad\longrightarrow\qquad
\text{子域对齐：}\frac1C\sum_{c=1}^{C}\mathcal D(P_s^c,P_t^c).
\]
\begin{columns}[T,onlytextwidth]
\column{.48\textwidth}
\begin{block}{全局统计混合了什么？}
整体协方差同时包含类内变化和类间中心结构。\par
不同类别的域偏移可能相互掩盖。
\end{block}
\column{.48\textwidth}
\begin{block}{子域对齐匹配什么？}
源域类别 $c$ 对应目标域类别 $c$，分别建立统计量并对齐。\par
目标域使用软类别权重，而非真实标签。
\end{block}
\end{columns}
\[
\mathcal L=\CE+\lambda\,\mathbb E_c[\mathcal D(P_s^c,P_t^c)].
\]
\src{\href{https://arxiv.org/pdf/2106.09388}{DSAN，第 III 节}：按类别划分相关子域。}
\end{frame}

\begin{frame}[t]{参考 DSAN 实现 CORAL 子域对齐}
\small
借用 DSAN 的类别软权重，将其统计目标替换为\textbf{逐类协方差差异}：
\[
p_{jc}=\sg[\operatorname{softmax}(z_j^t)_c],\quad
a^s_{ic}=\frac{\mathbf1[y_i^s=c]}{\sum_r\mathbf1[y_r^s=c]},\quad
a^t_{jc}=\frac{p_{jc}}{\sum_rp_{rc}}.
\]
\[
\mu_d^c=\sum_i a^d_{ic}z_i^d,\qquad
C_d^c=\frac{\sum_i a^d_{ic}(z_i^d-\mu_d^c)(z_i^d-\mu_d^c)^\top}{1-\sum_i(a^d_{ic})^2}.
\]
\[
\boxed{\LC=\frac1{|\mathcal V|}\sum_{c\in\mathcal V}\frac{\norm{C_s^c-C_t^c}_F^2}{4D^2}}.
\]
\begin{itemize}
  \item AlexNet fc8 的 31 维 logits；目标权重 detach，特征路径反传。
  \item 有效类要求：源支持 $\ge2$、目标 top-1 支持 $\ge2$、目标 ESS $\ge2$。
\end{itemize}
\src{src/core/subdomain\_coral.py；类别权重参考 DSAN，协方差目标为本实验自定义。}
\end{frame}

\begin{frame}[t]{CORAL 子域对齐：未产生正向效果}
\small
统一初始化、batch=128；适应组前十轮均使用 $0.5\G$，第 11 轮切换。
\vspace{2mm}
\begin{tabularx}{\linewidth}{@{}Ycr@{}}
\toprule
配置 & 第 11--50 轮适应项 & 第 50 轮正确率 \\
\midrule
\SubRows
\bottomrule
\end{tabularx}
\vspace{3mm}
\begin{itemize}
  \item local 相对 global \textbf{下降 11.70 pp}。
  \item mixed 在保留全局系数 0.25 的情况下，仍比 global\_quarter \textbf{低 4.53 pp}。
\end{itemize}
\takeaway{更细的类别划分没有改善分类，直接沿用类内协方差目标反而损害正确率。}
\src{reports/subdomain\_coral/lmmd/summary.csv；均为固定第 50 轮结果。}
\end{frame}
